\documentclass[12pt]{article}
\usepackage[utf8]{inputenc}
\usepackage[T1]{fontenc}
\usepackage{hyperref}
\title{Sabaic Language Resources for K-12 Curriculum Development}
\author{Compiled by Claude Sonnet 5.5 (claude-sonnet-5-5) for Andr\'as Kornai}
\date{}
\begin{document}
\maketitle

\section{Introduction}
Sabaic (Sabaean; ISO 639-3 xsa) is an extinct Ancient South Arabian (Old South Arabian) Semitic language. It was spoken from about 1000 BC to the 6th century AD and is the language of most surviving South Arabian inscriptions. There are no native speakers, so any curriculum is a heritage and history one, not a language-immersion one.

\section{Primary Educational Resources}
No school textbooks or primers for Sabaic were found. Material for teaching would come from the scholarly resources below.

\section{Linguistic and Reference Materials}
Biella's \textit{Dictionary of Old South Arabic: Sabaean Dialect} (Harvard Semitic Studies 25, Scholars Press, 1982) is the standard printed dictionary: \url{https://glottolog.org/resource/reference/id/457791}. The DFG-funded Jena project is building a comprehensive Sabaic dictionary based on all known inscriptions: \url{https://gepris.dfg.de/gepris/projekt/210365926}. Its online dictionary is reported at sabaweb.uni-jena.de. Overview articles: \url{https://en.wikipedia.org/wiki/Sabaic} and \url{https://en.wikipedia.org/wiki/Ancient_South_Arabian}.

\section{Religious and Community Resources}
None found for a living community. Inscriptions are the main religious and historical record.

\section{Digital and Audio Resources}
The Corpus of Sabaic Inscriptions is a digital edition of about 3400 inscriptions, hosted in the DASI archive at the University of Pisa (reported in search; a direct URL was not seen). A CNR record is at \url{https://iris.cnr.it/handle/20.500.14243/374864}. There is a community subreddit at \url{https://reddit.birdcat.cafe/r/OldSouthArabian}. No audio was found, and none can exist for the ancient pronunciation.

\section{Curriculum Development Context}
A K-12 use would be enrichment in history, writing systems and Semitic linguistics, built around the South Arabian script, selected inscriptions and the Sabaean kingdom. The scholarly tools above are aimed at university level and would need adapting.
\end{document}